\section{Reproducibility and software}
\label{app:repro}

\paragraph{Encoder.} We load the encoder through \texttt{timm} from the Hugging Face hub by the
checkpoint string \texttt{hf\_hub:\allowbreak mwalmsley/\allowbreak zoobot-\allowbreak
encoder-\allowbreak greyscale-\allowbreak convnext\_nano}. The Zoobot package itself is not
installed and its version is not recorded; v2.0.1 is the release in which Zoobot's
documentation lists the greyscale encoders as added. The greyscale checkpoint has a
single-channel stem, so the model must be built with one input channel. Our unfrozen setting
trains timm \texttt{stages[3]}, \texttt{stages[2]} and \texttt{stages[1]}, the encoder's own
fine-tuning convention for $n_{\mathrm{blocks}} = 3$; the final layer normalisation is never
trained; the benchmark states only that it fine-tunes the last three
blocks~\citep{lines2025simstosky}.

\paragraph{Library versions.} We did not record the \texttt{timm} version used in training or
pin it in a requirements file. The notebook for the unfrozen FreeMatch arm installs
\texttt{timm} 1.0.26, but its outputs do not record the installed version, and the notebooks
for the other unfrozen arms install none, so they used the hosted environment's version; the
arms may therefore have run under different \texttt{timm} versions. Comments in our code name
1.0.26 and 1.0.28 as checked, and the per-stage parameter counts are a regression test, so a
release that renumbered the stages would fail a test. We installed GalSim unpinned in the WSL2
conda environment used for simulation and did not record its version, so the catalogue's default selection cannot be tied to
a known release. Our input stretch reproduces the \texttt{fits\_to\_jpg\_like\_scaling}
function of the \texttt{galaxy-datasets} package from its source code; we did not record the
package version or commit.

\paragraph{Compute venue and precision.} Every unfrozen arm trained in \venueTrainPrecision{}
and was scored in \venueEvalPrecision{}. The \venueGpuArms{} runs record a \venueGpu{}; the
FreeMatch runs record only a CUDA device, but their paths place them on Kaggle, and the
notebook settings we ship (\texttt{KAGGLE\_FREEMATCH\_PERFOLD\_CELLS.md}) specify a T4. We
measured the effect
of the venue once, on the frozen baseline with weak augmentation at half precision: the same
configuration scored \augprecCellC{} locally and \augprecCellCKaggle{} on Kaggle, a difference
of \augprecVenueDiff{} with interval \augprecVenueCI{} (Appendix~\ref{sec:augprec}). A second
check did not reproduce a venue effect: the same unfrozen CORAL weights scored in two separately
stored scorings gave identical TPR@$10^{-3}$ at all \venueCoralCheckpoints{} shared checkpoints
and AUROC agreeing to \venueCoralMaxAuroc{}.

\paragraph{The 1:10 runs.} The diluted DANN, BNM and unfrozen CORAL runs record a \venueGpu{};
the FreeMatch runs, frozen and unfrozen, record only a CUDA device; and the CORAL frozen stage
ran in full precision on CPU. Before FreeMatch's unfrozen runs, the pre-run check that part 3 of
the pre-registration requires reproduced the pattern of FreeMatch's early skipped updates in
its original mode, but not their exact count; we judged the match sufficient and proceeded,
since the rule does not say how closely the known-null must match. Before computing anything new,
the analysis reproduced this paper's balanced per-fold TPRs and its \dlReproNWord{} contrasts
against the control exactly, at four decimals.

\paragraph{Skipped updates.} In half precision the gradient scaler skips any step whose
gradients overflow. On folds 1--4 the balanced runs skipped \skipLamList{} steps (the control),
\skipDannList{} (DANN) and \skipCoralList{} (CORAL), and FreeMatch \skipFmEach{} on every fold;
the BNM runs did not record skips. Table~\ref{tab:dl-bookkeeping} gives the 1:10 runs'.

\paragraph{Infrastructure failures.} Three failures cost us time. Continuing training from a
weights-only checkpoint silently reset the training state; Appendix~\ref{sec:appendix-restart}
gives the cost and the repair. A gradient-norm measurement taken on CPU in fp32 moved by
$+\rvGradShiftPct{}$\% when re-measured on a T4 in fp16, enough to shift the weight grid derived
from it, so we re-derived the grid on the training hardware. An audit that found no
hard-coded write paths in our scripts was accurate when made, but a script written after it
wrote to one; we found it later and gave it an output flag and an overwrite guard.

\paragraph{Hash enrolment.} Two mechanisms implement our hash protection: discovery of every
file whose digest our run logs record (in the release, the digests in
\texttt{public\_inputs/\allowbreak record\_extract.json}), and a separate inventory that the
analysis scripts verify at
start and end. The checkpoint-window declaration of Appendix~\ref{sec:failed} is hashed but was
never enrolled in that inventory, so it is neither re-checked automatically nor read-only.
Its digest currently matches the file on disk, and Table~\ref{tab:hashes6} re-derives it at build
time and marks it as the one uncovered row. We did not enrol it after the fact, because that
would date its protection to the writing of this paper. The measuring run behind one hashed
gradient-norm grid was never downloaded, and the per-checkpoint files of two unfrozen
development runs on fold 0 were lost (Section~\ref{sec:negative}).

\paragraph{Script provenance.} The loss module that our four $\lambda = 0$ confirmatory folds ran
had since been overwritten in our working copy. We recovered it from the archived training
bundle, and its digest matches the one those runs recorded. Through line 1440 it is identical to
the version our DANN folds ran; the remainder is a BNM section that the $\lambda = 0$ training
path does not import, as the pre-registration's own audit, recovered with its recorded digest,
also documents. Four script versions used only by fold-0 development runs remain unrecovered;
they affect development-fold values only. The version of the unfrozen FreeMatch runner that
ran our balanced FreeMatch fold 1 is also unrecovered: its run record lacks \fmFoldOneMissingWord{} keys that
today's runner always writes and records the exact step-zero tolerance that preceded our gate
change (Appendix~\ref{sec:appendix-restart}), so the released runner is a later version.
For \dlLostArmsWord{} of the 1:10 arms (DANN, BNM and the CORAL and FreeMatch frozen
stages), the exact versions of \dlLostFilesWord{} helper scripts that ran
(\texttt{dilute10\_common.py}, \texttt{dilute10\_batch\_log.py} and the arm's wrapper) are lost;
the release has their later versions. Their recorded digests are in the run files. For
Baseline A, FreeMatch and CORAL's balanced frozen stage, the runs recorded no per-module digests,
and BNM's balanced runs recorded only the files its pre-registration hashes.

\paragraph{Where things are.} The contrast set and its registration marks are
\texttt{CONTRAST\_ROWS} and \texttt{CONTRAST\_PREREG} in \texttt{paper/scripts/\_common.py};
the four registered contrasts are in \texttt{runs/lambda0\_r6\_prereg.json}; the test-pool
exclusion rule is \texttt{audit/decontam/rule.json}; and the post hoc analyses of
Sections~\ref{sec:resolution}, \ref{sec:findmiss} and Appendices~\ref{app:bright}
and~\ref{app:breakdowns} are in \texttt{audit/review\_r1/}, \texttt{audit/review\_r1b/} and
\texttt{audit/rewrite\_v5/}; the
scripts that generate the tables, figures and number macros are in \texttt{paper/scripts/}.
The 1:10 experiment's contrasts and their mark are \texttt{DILUTE10\_CONTRASTS} and
\texttt{DILUTE10\_PREREG} in \texttt{paper/scripts/\_dilute10.py}; its three-part
pre-registration is
\texttt{runs/\allowbreak dilute10\_prereg\_\allowbreak part\{1,2,3\}.json}, its pool
\texttt{data/\allowbreak splits/\allowbreak dilute10\_pool.json}, its analysis
\texttt{runs/\allowbreak dilute10\_analysis/}, the frozen-stage head norms
\texttt{runs/\allowbreak dilute10\_head\_norms/}, and the memorisation diagnostic
\texttt{runs/\allowbreak dann\_lambda0\_holdout\_gap/}. The three files from the archived
training bundle, the loss module of the $\lambda = 0$ confirmatory folds, the
pre-registration's audit and the pre-registration itself, are
\texttt{evidence/\allowbreak kaggle\_bundle\_v14/\allowbreak da\_losses.py},
\texttt{evidence/\allowbreak kaggle\_bundle\_v14/\allowbreak lambda0\_r6\_audit.json} and
\texttt{evidence/\allowbreak kaggle\_bundle\_v14/\allowbreak lambda0\_r6\_prereg.json}.
